\begin{helper}
Let \(I\subseteq[n]\) have
\[
|I|=\noderef{TwoBiteNaturalIndependenceScale}(1+\varepsilon,n).
\]
Assume the fiber-degree event
\(\noderef{FiberAndDegreeEvent}(\omega,\varepsilon_2)\), the eventual
comparison package
\(\noderef{ParameterHierarchyEventualComparisons}
(\varepsilon,\varepsilon_1,\varepsilon_2,n_0)\), \(n>n_0\), and the rounded
sample-parameter identities
\[
m=\noderef{TwoBiteNaturalReducedVertexCount}(n),\qquad
p=\frac12\sqrt{\frac{\log n}{n}}.
\]
If a base vertex \(x\) is huge for \(I\), then \(x\) is
\(\noderef{TwoBiteStageRevealedVertex}(\omega,\varepsilon,I,x)\).
\end{helper}
\begin{proof}
Write \(L=\log n\),
\[
t_1=\noderef{TwoBiteHugeCutoff}(n),\qquad
t_2=\noderef{TwoBiteLargeCutoff}(\varepsilon,n).
\]
The hugeness assumption is the conjunction defining
\(\noderef{TwoBiteHugeClass}(\omega,I,x)\).  Its first projection gives
\[
t_1<|\noderef{TwoBiteX}(\omega,I,x)|, \tag{H1}
\]
and its second projection says
\[
|\noderef{TwoBiteX}(\omega,I,x)|\le |I|. \tag{H2}
\]
The second inequality is also consistent with the definition of
\(\noderef{TwoBiteX}\), which is the filter of \(I\) consisting of vertices
lying in the lifted neighborhood of \(x\).

First suppose \(x\notin\pi_R(I)\cup\pi_B(I)\), equivalently
\(\neg\noderef{TwoBiteProjectionContains}(\omega,I,x)\).  Then the first
disjunct in the definition of
\(\noderef{TwoBiteStageRevealedVertex}\) holds, so \(x\) is staged-revealed.

It remains to consider the projection-contained case.  We prove the
contrapositive of the paper's inclusion
\[
(\pi_R(I)\cup\pi_B(I))\cap H_I\subseteq F_2.
\]
Assume
\[
\noderef{TwoBiteProjectionContains}(\omega,I,x)
\quad\text{and}\quad
\neg\noderef{TwoBiteStageF2}(\omega,\varepsilon,I,x).
\tag{P}
\]
Unfolding \(\noderef{TwoBiteStageF2}\), the statement
\(\noderef{TwoBiteStageF2}(\omega,\varepsilon,I,x)\) is the conjunction of
\(\noderef{TwoBiteProjectionContains}(\omega,I,x)\) and the strict real
inequality
\[
\frac{t_2}{L}<
\left|
\{y:\noderef{TwoBiteStageF1}(\omega,I,y)
      \text{ and }\noderef{TwoBiteBaseAdj}(\omega,x,y)\}
\right|.
\]
The first conjunct is already true by (P).  Hence, if the displayed strict
inequality also held, the two conjuncts would prove
\(\noderef{TwoBiteStageF2}(\omega,\varepsilon,I,x)\), contradicting (P).
Therefore the strict inequality is false, and totality of the real order gives
the weak bound on the same-color \(F_1\)-neighbor count:
\[
\left|
\{y:\noderef{TwoBiteStageF1}(\omega,I,y)
      \text{ and }\noderef{TwoBiteBaseAdj}(\omega,x,y)\}
\right|
\le \frac{t_2}{L}. \tag{N}
\]

We next split the finite set
\[
X_x(I)=\noderef{TwoBiteX}(\omega,I,x)
\]
according to the same-color projection of its elements.  If \(x=\mathrm{red}
(r)\), define the same-color projection of \(u\in I\) to be
\(\pi_R^\omega(u)\); if \(x=\mathrm{blue}(b)\), define it to be
\(\pi_B^\omega(u)\).  In either color case, membership
\(u\in X_x(I)\) unfolds to two facts: \(u\in I\), and the same-color projected
base vertex \(y\) is adjacent to \(x\) in the corresponding base graph.  Thus
\(\noderef{TwoBiteBaseAdj}(\omega,x,y)\) holds.  Since \(u\in I\) projects to
\(y\), the definition of \(\noderef{TwoBiteProjectionContains}\) also gives
\[
\noderef{TwoBiteProjectionContains}(\omega,I,y). \tag{C}
\]
Let \(A\) be the part of \(X_x(I)\) for which this same-color projected base
vertex \(y\) lies in \(\noderef{TwoBiteStageF1}(\omega,I,\cdot)\), and let
\(B\) be the complementary part.  These are complementary filters of the same
finite set, hence they are disjoint, cover \(X_x(I)\), and satisfy
\[
|X_x(I)|=|A|+|B|. \tag{D}
\]

For \(A\), every \(u\in A\) lies in the fiber over a base vertex \(y\) with
\[
\noderef{TwoBiteStageF1}(\omega,I,y)
\quad\text{and}\quad
\noderef{TwoBiteBaseAdj}(\omega,x,y).
\]
The number of such base vertices is bounded by (N).  For each such \(y\), the
appropriate red or blue fiber-size conjunct of
\(\noderef{FiberAndDegreeEvent}(\omega,\varepsilon_2)\) gives
\[
\noderef{WithinRelativeError}
  \left(\varepsilon_2,L^2,|F(y)|\right).
\]
Using \(0\le\varepsilon_2\) from
\(\noderef{ParameterHierarchyEventualComparisons}\) and \(L^2\ge0\), this
relative-error inequality implies the uniform upper bound
\[
|F(y)|\le (1+\varepsilon_2)L^2
\]
for every same-color fiber, in both the red and blue cases.  The finite
fiber-counting estimate
\[
|A|\le
\sum_{\substack{y:\ \noderef{TwoBiteStageF1}(\omega,I,y)\\
        \noderef{TwoBiteBaseAdj}(\omega,x,y)}} |F(y)|
\]
then gives
\[
|A|\le (1+\varepsilon_2)L^2\,\frac{t_2}{L}. \tag{A}
\]

For \(B\), the same-color projected base vertex \(y\) is still adjacent to
\(x\), but it is not in \(\noderef{TwoBiteStageF1}(\omega,I,\cdot)\).  The
base-degree conjunct of
\(\noderef{FiberAndDegreeEvent}(\omega,\varepsilon_2)\), using the red graph
when \(x\) is red and the blue graph when \(x\) is blue, gives the same-color
degree bound
\[
|\{y:\noderef{TwoBiteBaseAdj}(\omega,x,y)\}|\le (1+\varepsilon_2)pm.
\tag{G}
\]
Now fix a base vertex \(y\) that contributes at least one element to \(B\).
By (C), \(y\) is projection-contained.  Since \(y\notin
\noderef{TwoBiteStageF1}(\omega,I,\cdot)\), unfolding
\(\noderef{TwoBiteStageF1}\) shows that the strict inequality
\[
L<|F(y)\cap I|
\]
cannot hold; otherwise it would combine with projection-containedness to put
\(y\) in \(F_1\).  Hence
\[
|F(y)\cap I|\le L.
\]
Every element of \(B\) lies in \(F(y)\cap I\) for one of the same-color
neighbors \(y\) of \(x\).  Summing the preceding bound over all such neighbors
and then using (G) gives
\[
|B|\le L(1+\varepsilon_2)pm. \tag{B}
\]

Adding the bounds for \(A\) and \(B\), and using the disjoint decomposition
(D), gives
\[
|X_x(I)|
\le
(1+\varepsilon_2)L^2\,\frac{t_2}{L}
{}+
L(1+\varepsilon_2)pm. \tag{T}
\]
The rounded sample-parameter identities identify the actual \(m\) and \(p\)
in (T) with the \(m\) and \(p\) used in
\(\noderef{ParameterHierarchyEventualComparisons}\); the active scale
hypothesis similarly identifies \(|I|\) with the rounded
\(\noderef{TwoBiteNaturalIndependenceScale}(1+\varepsilon,n)\), the scale used
when the hierarchy package is instantiated at this \(n\).  The exact
comparison conjunct of that package says
\[
(1+\varepsilon_2)L^2\,\frac{t_2}{L}
{}+
L(1+\varepsilon_2)pm
<t_1. \tag{H}
\]
Combining (T) and (H) gives \(|X_x(I)|<t_1\), contradicting the huge lower
bound (H1).  Therefore the assumptions in (P) are inconsistent, so a
projection-contained huge vertex must lie in
\(F_2\).

In the projection-contained case, the \(F_2\) disjunct of
\(\noderef{TwoBiteStageRevealedVertex}\) holds.  Together with the
outside-projection case above, this proves the staged-revealed conclusion for
every huge vertex.
\end{proof}
